% Chapter 9 -- who could be interested in my profile. Update with background/* (screen.py, scheduler.py,
% probes.py, texts.py), suggest/* (brief.py, prompt.py, runner.py, engine.py, verify.py, texts.py),
% gui/interest_flow.py, gui/window_interest.py, gui/background_settings.py, gui/suggestions_view.py.
\chapter{Who could be interested in my profile}\label{ch:interest}

Tab 2 starts from a place and a radius. This chapter starts from the \emph{profile}: given the CV
and the career directions you want, which companies, known or not yet known, could want it, and
which deserve a full analysis first? The button \ui{Who could be interested in my profile?} (tab~1,
and \ui{Search $>$ Who could be interested in my profile?\ldots}) opens one non-modal window with four
boxes, \emph{1. Where and what} (an area: a town and a radius of 1--300\,km, a province, a region or
the whole country; the career directions of chapter~\ref{ch:affinity}; the four company types),
\emph{2. Companies already known, in order}, \emph{3. Ask your assistant} and \emph{4. Deep analysis}.
Closing the window cancels what runs in it; slow work shows in the shared progress bar.

\section{The fast pre-screen}\label{sec:prescreen}
It starts by itself when the window opens with a CV, uses \emph{no network} and takes well under a
second per thousand companies. It orders every company the app already knows, without hiding any:
\begin{itemize}
  \item \textbf{Candidates}: your saved cards, the regional register for the chosen area (up to 400,
        chapter~\ref{ch:register}) and the local map index if it is already on disk (up to 1\,500);
        duplicates are joined.
  \item \textbf{Score}, for each enabled direction: $0.45\cdot$terms $+\,0.25\cdot$type fit $+\,0.30\cdot$rank.
        \emph{Terms} counts role and skill phrases of the direction found in a short text built from the
        name, description, sector, products and register or map tags; \emph{type fit} is how well the
        guessed company type (the knowledge table of section~\ref{sec:knowledge}) suits the direction; \emph{rank}
        is the free pre-crawl ranking against the CV (chapter~\ref{ch:discovery}). The best direction wins and
        the score is multiplied by your company-type weight; a type set to 0 is kept but ordered last.
  \item \textbf{Table}: Company, Source, Place, \emph{Why (quick estimate)}, Analysis state (\emph{not
        analysed yet}, \emph{queued}, \emph{analysing}, \emph{analysed}) and the four affinity columns once a
        result exists. A direction guessed from the name and the tags alone (no matching word, no known company type) is
        shown as ``maybe \ldots (to be confirmed when the site is read)''. It is a quick estimate for ordering, \emph{not} an affinity; double-click opens the
        card, and the right-click menu offers \ui{Analyse selected now} or \ui{Analyse when the PC is idle}.
\end{itemize}

\section{Assistant suggestions}\label{sec:suggest}
For companies the app does not know, \ui{Suggest companies with the assistant} asks \emph{your own}
Claude Code or Codex (Settings, \ui{Assistant and MCP}; it uses your plan, not a ReverseHire service) for
up to 10 real companies in the area, best fit first; \ui{Show 10 more} continues the same conversation. The
same view is under the companies list of the job search.

\textbf{The anonymised brief.} Before the first request a window shows \emph{exactly} what will be sent, and
\ui{What is sent} reopens it. In English, one fact per line: up to 6 roles, seniority and years (to the half
year), up to 8 skill families with approximate years, up to 6 experiences as ``title, years, undisclosed
employer'', sectors, up to 5 languages, travel willingness, the wanted directions and company types, and the
area as words. Not sent: name, contacts, address, home town, employer names, dates, any CV text, the
coordinates of the area. The text is scrubbed against the CV's own names and addresses, and if anything still
leaks the personal parts are emptied. The names of companies you already know in the area go along as an exclusion list: at most 60 names (the
companies already suggested first, then your saved ones), \emph{shown in the same window}: a line gives their
number and \ui{Show the names} lists them. \ui{Show 10 more} adds the names already given in the conversation, never
more than 300 in all.

\textbf{A locked-down CLI.} The CLI runs in an empty folder, with the prompt on standard input. With Claude
Code every built-in tool is off, an empty MCP configuration is forced, nothing is asked and at most 2 turns
are allowed; with Codex the sandbox is read-only, approvals are never requested, MCP servers and shell, browser
and other features are off. Your own settings, hooks and rules are ignored, and any tool call other than the
optional web search stops the run. \textbf{Web search} is off by default (Settings, \ui{Idle analysis}):
the app verifies every company anyway, and search results are untrusted text. The reply must be one JSON
object, \code{\{"companies": [\ldots], "more\_available": \ldots\}}, each item with name, municipality,
province, country, website, company type, direction, one-sentence \emph{why} and a confidence;
the parser tolerates code fences and a bare list, and drops bad items one by one (a website must be a plain
public http(s) address).

\textbf{Nothing the assistant says is a fact until it is checked}, by the same polite fetcher as everywhere
(robots.txt, delays, a budget of 8 requests per company and 80 per batch):
\begin{enumerate}[leftmargin=1.5em]
  \item the \emph{website} must be reachable, not a parked domain, and carry the company's name in its title,
        site name, structured data or legal block (another company's legal name there is a contradiction);
  \item the \emph{local map index} (no network) is searched for the same company within 12\,km, and the
        \emph{regional register} (no network) for a company of the same name in the province;
  \item \emph{identity}: a VAT number printed on the site must check out and be confirmed by VIES with a
        matching name; a company tax code is looked up in GLEIF;
  \item \emph{Wikidata} is asked, by name, only when the above did not settle it; one item with a nearly identical label
        and an official website is accepted;
  \item the \emph{place} is compared with the requested area (a confirmed position outside is a contradiction).
\end{enumerate}

{\small
\begin{tabularx}{\textwidth}{@{}l Y@{}}
\toprule
\hdr \thead{Status} & \thead{Meaning} \\
\midrule
\textcolor{othGreen}{\textbf{Verified}}  & Site reachable with the name on it, plus identity confirmed, or the same company on the map, or a Wikidata item with the same website. \\
\textcolor{othBlue}{\textbf{Partly verified}} & Site reachable with the name on it, nothing else confirms it; or no readable site, but the regional register lists the company in the area. \\
\textcolor{othGray}{\textbf{Not verified}} & No reachable site carrying the name: shown greyed, can never be added; absence of evidence is not ``does not exist''. \\
\textcolor{othRed}{\textbf{Rejected}}     & Contradicted (the site belongs to another company, a VAT number to another, the place is outside the area); hidden until \ui{Show rejected} is ticked. \\
\bottomrule
\end{tabularx}
}

\vspace{3pt}
The table shows the status, the company, place, the website the app read, the checks as
confirmed/contradicted lines and, apart, \emph{``Assistant's opinion (not verified)''}. \ui{Add to my companies}
(only for verified or partly verified) creates a card tagged ``suggested by assistant'' with the checks in its
notes. A request is capped at 240\,s (900\,s with verification); a notice under the table shows the tokens or cost the
CLI reported, if any. If your plan has reached its usage limit the window says so, with the reset time when the
assistant gave one. There is no daily limit on requests.

\section{Deep analysis and the idle scheduler}
\ui{Analyse the best now} (box 4, 1--50 companies, default 10) builds the full card of the best rows one after
the other (chapter~\ref{ch:card}; the paid provider is never used), then computes their affinity, filling the
columns; the rest are \emph{queued}: at most 200, the best first, and the first time the app asks whether to do so
(it tells you the idle time, mains power and daily budget that apply) and remembers the answer. A queued company is analysed by the \textbf{idle scheduler}, which works
only when you are not using the PC:

{\small
\begin{tabularx}{\textwidth}{@{}l Y@{}}
\toprule
\hdr \thead{Condition} & \thead{Rule (defaults; Settings, \ui{Idle analysis})} \\
\midrule
Idle time          & No input for at least 5 minutes (1--240). Read from the system (Windows: last input; Linux X11: \code{xprintidle}; macOS: HID idle time); if unknown, the app must have been in the background for at least 10 minutes. Input inside the app pauses at once. \\
Nothing else running & No company search, job search, card build, suggestion request or work of the interest window. \\
Power, CPU, network & On mains power (setting, on), CPU load below 50\,\%, network up; ``unknown'' counts as allowed. \\
Budget             & At most 600 web requests a day (25--20\,000), 60 a day per site, one job at a time. \\
\bottomrule
\end{tabularx}
}

\vspace{3pt}
When a condition fails, the running job is cancelled at its next wait, goes back to the queue \emph{without
penalty}, keeps the steps already saved and resumes by itself. A failure is retried after 5\,min, doubling up to 6\,h,
and given up after 5 attempts. The best pre-screen score goes first. The status line shows ``Idle analysis: 4
done, 12 queued, waiting'' with the reason; \ui{Tools $>$ Pause the idle analysis} (or the button in the bar)
pauses it, and the pause survives a restart. The queue and the day's counters are kept in \code{idle\_analysis.json}
and read again at start-up; nothing is analysed at start-up unless the conditions hold.

\begin{infobox}
Only the requests of the site fetcher are counted against the daily budget, and after each job, so a job can
overshoot it slightly. The ranking of this chapter never uses your name, address or CV text: it uses the profile.
\end{infobox}
